% Part: proof-theory
% Chapter: natural-deduction
% Section: rules-N2

\documentclass[../../../include/open-logic-section]{subfiles}

\begin{document}

\begin{table}
\begin{tabular}{@{}c@{}c@{}}
\multicolumn{2}{l}{\textbf{अभिगृहीत:}}\\
\multicolumn{2}{c}{$\Gamma \Sequent x:!A$, यदि $x:!A \in \Gamma$}\\[3ex]
\multicolumn{2}{l}{\textbf{अनुमान-नियम:}}\\[3ex]
\multicolumn{2}{c}{\Axiom$x: !A, \Gamma \fCenter N:!B$
\RightLabel{\Intro{\lif}}
\UnaryInf$\Gamma \fCenter \lambd[\typeof{x}{!A}][N] : !A \lif !B $
\DisplayProof}
\\[4ex]
\multicolumn{2}{c}{\Axiom$\Gamma \fCenter N:!A \lif !B$
\Axiom$\Gamma \fCenter M:!A$
\RightLabel{\Elim{\lif}}
\BinaryInf$\Gamma \fCenter (NM) : !B$
\DisplayProof}
\\[4ex]
\multicolumn{2}{c}{\Axiom$\Gamma \fCenter N_1 : !A$
\Axiom$\Gamma \fCenter N_2 : !B$
\RightLabel{\Intro{\land}}
\BinaryInf$\Gamma \fCenter \pair{N_1}{N_2} : !A \land !B $
\DisplayProof}
\\[4ex]
\Axiom$\Gamma \fCenter N: !A \land !B$
\RightLabel{\Elim{\land}[1]}
\UnaryInf$\Gamma \fCenter \proj{1}{N} : !A$
\DisplayProof
&
\Axiom$\Gamma \fCenter N:!A \land !B$
\RightLabel{\Elim{\land}[2]}
\UnaryInf$\Gamma \fCenter \proj{2}{N} : !B$
\DisplayProof\\[3ex]
\\[4ex]
\Axiom$\Gamma \fCenter N: !A$
\RightLabel{\Intro{\lor}[1]}
\UnaryInf$\Gamma \fCenter \inj[!B]{1}{N} : !A \lor !B$
\DisplayProof
&
\Axiom$ \Gamma \fCenter N : !A$
\RightLabel{\Intro{\lor}[2]}
\UnaryInf$\Gamma \fCenter \inj[!B]{2}{N} : !B \lor !A$
\DisplayProof\\[3ex]
\multicolumn{2}{c}{
\begin{tabular}[b]{@{}l@{}}
\Axiom$\Gamma \fCenter M : !A \lor !B$
\Axiom$x:!A, \Gamma \fCenter N_1 : !C$
\Axiom$y:!B, \Gamma \fCenter N_2: !C$
\RightLabel{\Elim{\lor}}
\TrinaryInf$\Gamma \fCenter
\dcase{M}{x}{N_1}{y}{N_2} : !C$
\DisplayProof
\end{tabular}
}
\\[4ex]
\multicolumn{2}{c}{\Axiom$\Gamma \fCenter N: \lfalse$
\RightLabel{\FalseInt}
\UnaryInf$\Gamma \fCenter \abort{!A}{N} : !A$
\DisplayProof}
\end{tabular}
\caption{प्रकार-निर्धारण प्रणाली के रूप में प्रयुक्त पद-चिह्नित प्रतिज्ञप्तिपरक
अंतःप्रज्ञात्मक प्राकृतिक निगमन प्रणाली \Log{tN3ip} के नियम।}
\ollabel{tab:tN2ip}
\end{table}

\end{document}
